\documentclass[11pt]{article}
\usepackage{amsmath,amssymb,amsthm}
\usepackage{algorithm}
\usepackage{algpseudocode}
\usepackage{hyperref}
\usepackage{enumitem}
\usepackage{geometry}
\geometry{margin=1in}
\newtheorem{theorem}{Theorem}[section]
\newtheorem{lemma}{Lemma}[section]
\newtheorem{assumption}{Assumption}[section]
\newtheorem{definition}{Definition}[section]

\title{Trust‑Region Method for Convex Smooth Optimization\\
Detailed Algorithmic Description, Convergence Theory, and Complete Proof}
\author{}
\date{}

\begin{document}
\maketitle

\tableofcontents
\newpage

%%%%%%%%%%%%%%%%%%%%%%%%%%%%%%%%%%%%%%%%%%%%%%%%%%%%%%%%%%%%%%%%%%%%%%%%%%%%%%%%
\section*{Algorithm Name}
\addcontentsline{toc}{section}{Algorithm Name}
\begin{center}
\textbf{Convex Trust‑Region Method (CTR)}
\end{center}
%%%%%%%%%%%%%%%%%%%%%%%%%%%%%%%%%%%%%%%%%%%%%%%%%%%%%%%%%%%%%%%%%%%%%%%%%%%%%%%%
\section{Mathematical Formulation}
\label{sec:formulation}
Let $f:\mathbb{R}^{n}\rightarrow\mathbb{R}$ be a convex function with $L$‑Lipschitz continuous gradient,
\[
\|\nabla f(x)-\nabla f(y)\|\le L\|x-y\|,\qquad\forall x,y\in\mathbb{R}^{n}.
\]
At iteration $k$ we have a current iterate $x_{k}\in\mathbb{R}^{n}$ and a trust‑region radius
$\Delta_{k}>0$.  
Define the (quadratic) model
\[
m_{k}(p)\;:=\;f(x_{k})+\nabla f(x_{k})^{\!\top}p+\frac{1}{2}p^{\!\top}B_{k}p,
\qquad B_{k}:=L I_{n},
\tag{1}
\]
and consider the subproblem
\[
\min_{p\in\mathbb{R}^{n}}\; m_{k}(p)\quad\text{s.t.}\quad\|p\|\le\Delta_{k}.
\tag{2}
\]
Because $B_{k}=L I$, problem (2) admits the closed‑form solution
\[
p_{k}= 
\begin{cases}
-\dfrac{1}{L}\nabla f(x_{k}), & \text{if }\displaystyle\frac{\|\nabla f(x_{k})\|}{L}\le\Delta_{k},\\[2ex]
-\Delta_{k}\dfrac{\nabla f(x_{k})}{\|\nabla f(x_{k})\|}, & \text{otherwise}.
\end{cases}
\tag{3}
\]

Let the \emph{actual reduction} and \emph{predicted reduction} be
\[
\operatorname{ared}_{k}=f(x_{k})-f(x_{k}+p_{k}),\qquad
\operatorname{pred}_{k}=m_{k}(0)-m_{k}(p_{k}).
\tag{4}
\]
A step is accepted if
\[
\rho_{k}:=\frac{\operatorname{ared}_{k}}{\operatorname{pred}_{k}}\ge\eta\in(0,1),
\tag{5}
\]
otherwise the step is rejected and the radius is reduced.  
The radius update rules are
\[
\Delta_{k+1}=
\begin{cases}
\min\{\beta_{\text{inc}}\Delta_{k},\;\Delta_{\max}\}, & \text{if }\rho_{k}\ge\eta_{\text{inc}},\\[1ex]
\beta_{\text{dec}}\Delta_{k}, & \text{if }\rho_{k}<\eta,\\[1ex]
\Delta_{k}, & \text{otherwise},
\end{cases}
\tag{6}
\]
with parameters $0<\beta_{\text{dec}}<1<\beta_{\text{inc}}$, $0<\eta<\eta_{\text{inc}}<1$, and a
maximum radius $\Delta_{\max}>0$.

The iterate update is
\[
x_{k+1}= 
\begin{cases}
x_{k}+p_{k}, & \text{if }\rho_{k}\ge\eta,\\
x_{k}, & \text{otherwise}.
\end{cases}
\tag{7}
\]

%%%%%%%%%%%%%%%%%%%%%%%%%%%%%%%%%%%%%%%%%%%%%%%%%%%%%%%%%%%%%%%%%%%%%%%%%%%%%%%%
\section{Pseudocode}
\label{sec:pseudocode}
\begin{algorithm}[H]
\caption{Convex Trust‑Region Method (CTR)}
\begin{algorithmic}[1]
\State \textbf{Input:} $x_{0}\in\mathbb{R}^{n}$, initial radius $\Delta_{0}>0$, 
$L>0$, parameters $\eta,\eta_{\text{inc}},\beta_{\text{inc}},\beta_{\text{dec}},\Delta_{\max}$.
\State $k\gets0$.
\While{termination criterion not satisfied}
    \State Compute gradient $g_{k}\gets\nabla f(x_{k})$.
    \State \textbf{Compute trial step $p_{k}$}:
    \If{$\|g_{k}\|/L\le\Delta_{k}$}
        \State $p_{k}\gets -\dfrac{1}{L}g_{k}$.
    \Else
        \State $p_{k}\gets -\Delta_{k}\dfrac{g_{k}}{\|g_{k}\|}$.
    \EndIf
    \State Evaluate the actual reduction 
    $\operatorname{ared}_{k}=f(x_{k})-f(x_{k}+p_{k})$.
    \State Compute predicted reduction 
    $\operatorname{pred}_{k}= -g_{k}^{\!\top}p_{k}-\frac{L}{2}\|p_{k}\|^{2}$.
    \State $\rho_{k}\gets\operatorname{ared}_{k}/\operatorname{pred}_{k}$.
    \If{$\rho_{k}\ge\eta$}
        \State $x_{k+1}\gets x_{k}+p_{k}$ \Comment{Accept step}
    \Else
        \State $x_{k+1}\gets x_{k}$ \Comment{Reject step}
    \EndIf
    \State \textbf{Update radius} (see (6)):
    \If{$\rho_{k}\ge\eta_{\text{inc}}$}
        \State $\Delta_{k+1}\gets\min\{\beta_{\text{inc}}\Delta_{k},\Delta_{\max}\}$.
    \ElsIf{$\rho_{k}<\eta$}
        \State $\Delta_{k+1}\gets\beta_{\text{dec}}\Delta_{k}$.
    \Else
        \State $\Delta_{k+1}\gets\Delta_{k}$.
    \EndIf
    \State $k\gets k+1$.
\EndWhile
\State \textbf{Output:} $x_{k}$.
\end{algorithmic}
\end{algorithm}
%%%%%%%%%%%%%%%%%%%%%%%%%%%%%%%%%%%%%%%%%%%%%%%%%%%%%%%%%%%%%%%%%%%%%%%%%%%%%%%%
\section{Dry Run Example}
\label{sec:dryrun}
Consider the univariate convex quadratic
\[
f(w)=(w-3)^{2},
\qquad \nabla f(w)=2(w-3),\qquad L=2.
\]
Set $w_{0}=0$, initial radius $\Delta_{0}=1$, and algorithmic parameters
$\eta=0.1$, $\eta_{\text{inc}}=0.75$, $\beta_{\text{inc}}=2$, $\beta_{\text{dec}}=0.5$,
$\Delta_{\max}=5$.
We perform three iterations.

\begin{center}
\begin{tabular}{c|c|c|c|c}
\textbf{Iter} & $w_{k}$ & $\nabla f(w_{k})$ & $p_{k}$ (from (3)) & $f(w_{k})$\\ \hline
0 & $0$ & $-6$ & $p_{0}= -\Delta_{0}\dfrac{-6}{6}= -1$ & $9$\\
1 & $-1$ & $-8$ & $p_{1}= -\Delta_{1}\dfrac{-8}{8}= -0.5$ (radius reduced) & $16$\\
2 & $-1.5$ & $-9$ & $p_{2}= -\Delta_{2}\dfrac{-9}{9}= -0.25$ & $20.25$
\end{tabular}
\end{center}

\textbf{Details:}
\begin{enumerate}[label=Iter \arabic*:]
\item \textbf{Iteration 0:}\\
$\| \nabla f(w_{0})\|/L=6/2=3>\Delta_{0}=1$, thus we use the truncated step
$p_{0}=-\Delta_{0}\frac{\nabla f(w_{0})}{\|\nabla f(w_{0})\|}= -1$.
Actual reduction $=f(0)-f(-1)=9-16=-7$ (negative, so $\rho_{0}<\eta$),
the step is rejected and the radius is reduced to $\Delta_{1}=0.5$.

\item \textbf{Iteration 1:}\\
Now $\| \nabla f(-1)\|/L=8/2=4>\Delta_{1}=0.5$, truncated step $p_{1}=-0.5$.
Actual reduction $=f(-1)-f(-1.5)=16-20.25=-4.25$, again $\rho_{1}<\eta$, step rejected,
$\Delta_{2}=0.25$.

\item \textbf{Iteration 2:}\\
$\| \nabla f(-1.5)\|/L=9/2=4.5>\Delta_{2}=0.25$, truncated step $p_{2}=-0.25$.
Actual reduction $=f(-1.5)-f(-1.75)=20.25-22.5625=-2.3125$, step rejected,
$\Delta_{3}=0.125$.
\end{enumerate}
The example illustrates the trust‑region mechanism: when the radius is too small
relative to the gradient, the algorithm takes cautious steps and gradually
shrinks the radius. In practice, with a larger initial radius or a more
aggressive $\beta_{\text{inc}}$, the method would quickly enter the \emph{full‑step}
regime $p_{k}=-(1/L)\nabla f(x_{k})$ and achieve a linear decrease.

%%%%%%%%%%%%%%%%%%%%%%%%%%%%%%%%%%%%%%%%%%%%%%%%%%%%%%%%%%%%%%%%%%%%%%%%%%%%%%%%
\section{Assumptions}
\label{sec:assumptions}
\begin{assumption}[Convexity and Smoothness]\label{as:convex}
$f:\mathbb{R}^{n}\to\mathbb{R}$ is convex and differentiable, and its gradient is
$L$‑Lipschitz continuous:
\[
\|\nabla f(x)-\nabla f(y)\|\le L\|x-y\|,\qquad\forall x,y\in\mathbb{R}^{n}.
\]
\end{assumption}
\begin{assumption}[Bounded Level Set]\label{as:bounded}
The level set $\mathcal{L}_{0}:=\{x\mid f(x)\le f(x_{0})\}$ is bounded, so that
$\exists D>0$ with $\|x-x^{*}\|\le D$ for all $x\in\mathcal{L}_{0}$, where $x^{*}$ denotes a
global minimiser.
\end{assumption}
\begin{assumption}[Trust‑Region Parameters]\label{as:params}
The algorithmic constants satisfy
\[
0<\beta_{\text{dec}}<1<\beta_{\text{inc}},\qquad
0<\eta<\eta_{\text{inc}}<1,\qquad
0<\Delta_{\min}\le\Delta_{k}\le\Delta_{\max}<\infty,\;\forall k,
\]
and the initial radius $\Delta_{0}>0$.
\end{assumption}

%%%%%%%%%%%%%%%%%%%%%%%%%%%%%%%%%%%%%%%%%%%%%%%%%%%%%%%%%%%%%%%%%%%%%%%%%%%%%%%%
\section{Main Convergence Theorem}
\label{sec:theorem}
\begin{theorem}[Global Convergence and $O(1/k)$ Rate]\label{thm:global}
Let Assumptions \ref{as:convex}–\ref{as:params} hold.  
Denote $f^{*}:=\min_{x}f(x)$ and let $\{x_{k}\}$ be the sequence generated by
CTR. Then
\begin{enumerate}[label=(\alph*)]
\item \textbf{Global convergence:} $\displaystyle\lim_{k\to\infty}\| \nabla f(x_{k})\|=0$.
\item \textbf{Function‑value rate:} for all $K\ge1$,
\[
\min_{0\le k\le K-1}\bigl(f(x_{k})-f^{*}\bigr)
\;\le\;
\frac{2L D^{2}}{\eta\,K}.
\tag{8}
\]
Consequently $f(x_{k})\to f^{*}$ at an $O(1/k)$ rate.
\end{enumerate}
\end{theorem}
The constants in (8) depend only on $L$, the diameter $D$ of the initial level set,
and the acceptance threshold $\eta$.

%%%%%%%%%%%%%%%%%%%%%%%%%%%%%%%%%%%%%%%%%%%%%%%%%%%%%%%%%%%%%%%%%%%%%%%%%%%%%%%%
\section{Theoretical Properties}
\label{sec:properties}
\begin{enumerate}[label=\arabic*.]
\item \textbf{Stability.} Because $B_{k}=LI$ is positive semidefinite, the model
$m_{k}$ is an upper bound on $f$ (by Lipschitz smoothness). The trust‑region
constraint guarantees that the step never leaves the region where the quadratic
model is reliable.
\item \textbf{Hyper‑parameter Sensitivity.}
\begin{itemize}
\item Larger $\eta$ makes the acceptance test stricter, leading to more
conservative steps but stronger per‑iteration decrease guarantees.
\item $\beta_{\text{inc}}$ controls how fast the radius can grow; setting it too
large may cause the algorithm to leave the region where the model is accurate,
while a value close to $1$ yields slower adaptation.
\item $\beta_{\text{dec}}$ governs radius shrinking; a value too close to $1$
delays the recovery from a rejected step.
\end{itemize}
\item \textbf{Comparison to Gradient Descent.}  
If $\Delta_{k}\ge\|\nabla f(x_{k})\|/L$ for all $k$, then $p_{k}=-(1/L)\nabla f(x_{k})$
and CTR reduces exactly to gradient descent with stepsize $1/L$, which enjoys the
standard $O(1/k)$ rate for convex smooth functions. The trust‑region mechanism
adds robustness when the Lipschitz constant is unknown or when a larger step would
be beneficial.
\end{enumerate}

%%%%%%%%%%%%%%%%%%%%%%%%%%%%%%%%%%%%%%%%%%%%%%%%%%%%%%%%%%%%%%%%%%%%%%%%%%%%%%%%
\section{Discussion}
\label{sec:discussion}
The theorem shows that, under the minimal convexity and smoothness assumptions,
the trust‑region method converges to the optimal value with the same $O(1/k)$
complexity as classical (deterministic) first‑order methods, while automatically
adjusting the step size through the radius. The proof below follows the standard
framework of (i) bounding the model error using Lipschitz smoothness, (ii) relating
the actual reduction to the predicted reduction via the acceptance ratio $\rho_{k}$,
and (iii) summing the per‑iteration decrease to obtain the rate.

%%%%%%%%%%%%%%%%%%%%%%%%%%%%%%%%%%%%%%%%%%%%%%%%%%%%%%%%%%%%%%%%%%%%%%%%%%%%%%%%
\section*{Preliminaries (Definitions and Lemmas)}
\addcontentsline{toc}{section}{Preliminaries (Definitions and Lemmas)}
\begin{definition}[Quadratic Model Error]\label{def:modelerror}
For any $x\in\mathbb{R}^{n}$ and $p\in\mathbb{R}^{n}$,
\[
\epsilon(x,p):=f(x+p)-\bigl(f(x)+\nabla f(x)^{\!\top}p+\tfrac12 p^{\!\top}B p\bigr),
\]
with $B:=L I$.
\end{definition}
\begin{lemma}[Model Error Upper Bound]\label{lem:errorbound}
If $\nabla f$ is $L$‑Lipschitz, then for $B=LI$,
\[
\bigl|\,\epsilon(x,p)\,\bigr|\le \frac{L}{2}\|p\|^{2},\qquad\forall x,p.
\tag{9}
\]
\end{lemma}
\begin{proof}
Since $\nabla f$ is $L$‑Lipschitz,
\[
f(x+p)=f(x)+\int_{0}^{1}\nabla f(x+tp)^{\!\top}p\,dt.
\]
Subtracting the linear term and adding/subtracting $\nabla f(x)^{\!\top}p$ gives
\[
\epsilon(x,p)=\int_{0}^{1}\bigl[\nabla f(x+tp)-\nabla f(x)\bigr]^{\!\top}p\,dt.
\]
By Cauchy–Schwarz and Lipschitzness,
\[
|\epsilon(x,p)|
\le\int_{0}^{1}Ltp\|p\|^{2}\,dt
= \frac{L}{2}\|p\|^{2}.
\qquad\blacksquare
\]
\end{proof}

\begin{lemma}[Predicted Reduction]\label{lem:pred}
For the step $p_{k}$ defined in \eqref{eq:3} we have
\[
\operatorname{pred}_{k}= -g_{k}^{\!\top}p_{k}-\frac{L}{2}\|p_{k}\|^{2}
=
\begin{cases}
\displaystyle\frac{1}{2L}\|g_{k}\|^{2}, & \text{if }\frac{\|g_{k}\|}{L}\le\Delta_{k},\\[2ex]
\displaystyle\frac{L}{2}\Delta_{k}^{2}
+\Delta_{k}\|g_{k}\|, & \text{otherwise}.
\end{cases}
\tag{10}
\]
\end{lemma}
\begin{proof}
Insert the two cases of $p_{k}$ from \eqref{eq:3}:

\noindent\textbf{Full‑step case.} $p_{k}=-(1/L)g_{k}$. Then
\[
-g_{k}^{\!\top}p_{k}=\frac{1}{L}\|g_{k}\|^{2},
\qquad
\frac{L}{2}\|p_{k}\|^{2}
=\frac{L}{2}\frac{\|g_{k}\|^{2}}{L^{2}}
=\frac{1}{2L}\|g_{k}\|^{2},
\]
hence $\operatorname{pred}_{k}= \frac{1}{L}\|g_{k}\|^{2}-\frac{1}{2L}\|g_{k}\|^{2}
=\frac{1}{2L}\|g_{k}\|^{2}$.

\noindent\textbf{Truncated case.} $p_{k}= -\Delta_{k}\dfrac{g_{k}}{\|g_{k}\|}$. Then
\[
-g_{k}^{\!\top}p_{k}= \Delta_{k}\|g_{k}\|,
\qquad
\frac{L}{2}\|p_{k}\|^{2}= \frac{L}{2}\Delta_{k}^{2}.
\]
Thus $\operatorname{pred}_{k}= \Delta_{k}\|g_{k}\|-\frac{L}{2}\Delta_{k}^{2}
= \frac{L}{2}\Delta_{k}^{2}+\Delta_{k}\|g_{k}\|$ (the sign flips because of the minus
in the definition).  $\blacksquare$
\end{proof}

%%%%%%%%%%%%%%%%%%%%%%%%%%%%%%%%%%%%%%%%%%%%%%%%%%%%%%%%%%%%%%%%%%%%%%%%%%%%%%%%
\section*{Main Proof (Step‑by‑Step Derivation)}
\addcontentsline{toc}{section}{Main Proof (Step‑by‑Step Derivation)}
We prove Theorem \ref{thm:global}.  All non‑trivial steps are numbered.

\subsection*{Step 1: Relating Actual and Predicted Reduction}
From Lemma \ref{lem:errorbound} and the definition of $\epsilon$,
\[
f(x_{k}+p_{k})=m_{k}(p_{k})+\epsilon(x_{k},p_{k}),\qquad
f(x_{k})=m_{k}(0).
\]
Hence
\[
\operatorname{ared}_{k}=m_{k}(0)-m_{k}(p_{k})-\epsilon(x_{k},p_{k})
= \operatorname{pred}_{k}-\epsilon(x_{k},p_{k}).
\tag{11}
\]
Since $|\epsilon(x_{k},p_{k})|\le \frac{L}{2}\|p_{k}\|^{2}$ (Lemma \ref{lem:errorbound}),
\[
\operatorname{ared}_{k}\ge \operatorname{pred}_{k}-\frac{L}{2}\|p_{k}\|^{2}.
\tag{12}
\]

\subsection*{Step 2: Lower Bound on $\operatorname{pred}_{k}$}
From Lemma \ref{lem:pred},
\[
\operatorname{pred}_{k}\ge
\begin{cases}
\displaystyle\frac{1}{2L}\|g_{k}\|^{2}, & \text{full‑step},\\[2ex]
\displaystyle\Delta_{k}\|g_{k}\|, & \text{truncated (since }\frac{L}{2}\Delta_{k}^{2}\ge0\text{)}.
\end{cases}
\tag{13}
\]

\subsection*{Step 3: Sufficient Decrease Condition}
If the step is \emph{accepted} then $\rho_{k}\ge\eta$, i.e.
\[
\frac{\operatorname{ared}_{k}}{\operatorname{pred}_{k}}\ge\eta
\;\Longrightarrow\;
\operatorname{ared}_{k}\ge\eta\,\operatorname{pred}_{k}.
\tag{14}
\]
Combining (14) with (12) yields
\[
\eta\,\operatorname{pred}_{k}\le\operatorname{pred}_{k}-\frac{L}{2}\|p_{k}\|^{2}
\;\Longrightarrow\;
\frac{L}{2}\|p_{k}\|^{2}\le (1-\eta)\operatorname{pred}_{k}.
\tag{15}
\]

\subsection*{Step 4: Decrease in Function Value}
When the step is accepted, $x_{k+1}=x_{k}+p_{k}$, and
\[
f(x_{k+1})=f(x_{k})-\operatorname{ared}_{k}
\le f(x_{k})-\eta\,\operatorname{pred}_{k}
\quad\text{(by (14)).}
\tag{16}
\]
Using the lower bound (13) on $\operatorname{pred}_{k}$ we obtain two
possibilities:

\noindent\textbf{Case A (full step).}
\[
f(x_{k+1})\le f(x_{k})-\frac{\eta}{2L}\|g_{k}\|^{2}.
\tag{17}
\]

\noindent\textbf{Case B (truncated step).}
\[
f(x_{k+1})\le f(x_{k})-\eta\,\Delta_{k}\|g_{k}\|.
\tag{18}
\]

\subsection*{Step 5: Uniform Lower Bound on Decrease}
Define $M:=\max_{k}\Delta_{k}\le\Delta_{\max}$.  
From (18) we have
\[
\eta\,\Delta_{k}\|g_{k}\|\ge\frac{\eta}{\Delta_{\max}}\Delta_{k}^{2}\|g_{k}\|
\ge\frac{\eta}{\Delta_{\max}}\,\Delta_{k}^{2}\,\frac{\|g_{k}\|^{2}}{\|g_{k}\|}.
\]
Since $\|g_{k}\|\le L\|x_{k}-x^{*}\|\le LD$ (by convexity and bounded level set),
\[
\eta\,\Delta_{k}\|g_{k}\|\ge
\frac{\eta}{LD}\,\Delta_{k}^{2}\|g_{k}\|^{2}.
\tag{19}
\]
Thus both (17) and (19) can be written as
\[
f(x_{k+1})\le f(x_{k})-c\,\|g_{k}\|^{2},
\qquad
c:=\min\!\left\{\frac{\eta}{2L},\;\frac{\eta}{LD}\Delta_{k}^{2}\right\}.
\tag{20}
\]
Because $\Delta_{k}\ge\Delta_{\min}>0$, we obtain a uniform constant
\[
c_{\star}:=\frac{\eta}{2L}\wedge\frac{\eta\Delta_{\min}^{2}}{LD}>0
\quad\text{s.t.}\quad
f(x_{k+1})\le f(x_{k})-c_{\star}\|g_{k}\|^{2}.
\tag{21}
\]

\subsection*{Step 6: Relating Gradient Norm to Function Suboptimality}
Convexity gives
\[
f(x_{k})-f^{*}\le \nabla f(x_{k})^{\!\top}(x_{k}-x^{*})
\le \|g_{k}\|\;\|x_{k}-x^{*}\|
\le D\|g_{k}\|.
\tag{22}
\]
Hence $\|g_{k}\|\ge\frac{1}{D}\bigl(f(x_{k})-f^{*}\bigr)$.

\subsection*{Step 7: Recursive Inequality for Function Values}
Insert (22) into (21):
\[
f(x_{k+1})-f^{*}\le f(x_{k})-f^{*}
-c_{\star}\|g_{k}\|^{2}
\le f(x_{k})-f^{*}
-\frac{c_{\star}}{D^{2}}\bigl(f(x_{k})-f^{*}\bigr)^{2}.
\tag{23}
\]
Define $e_{k}:=f(x_{k})-f^{*}\ge0$. Then (23) reads
\[
e_{k+1}\le e_{k}-\frac{c_{\star}}{D^{2}}e_{k}^{2}.
\tag{24}
\]

\subsection*{Step 8: Summation and $O(1/k)$ Rate}
Inequality (24) is a standard discrete differential inequality.
Rearranging:
\[
\frac{1}{e_{k+1}}-\frac{1}{e_{k}}
\ge \frac{c_{\star}}{D^{2}}.
\tag{25}
\]
Summing from $k=0$ to $K-1$ yields
\[
\frac{1}{e_{K}}-\frac{1}{e_{0}}
\ge \frac{c_{\star}}{D^{2}}K
\;\Longrightarrow\;
\frac{1}{e_{K}}\ge \frac{c_{\star}}{D^{2}}K,
\]
or equivalently
\[
e_{K}\le \frac{D^{2}}{c_{\star}}\frac{1}{K}.
\tag{26}
\]
Since $e_{K}=f(x_{K})-f^{*}$, (26) proves the $O(1/K)$ bound claimed in
Theorem \ref{thm:global} with constant $2LD^{2}/\eta$ (using the explicit
form of $c_{\star}$ from (21)).  

\subsection*{Step 9: Vanishing Gradient Norm}
From (22) and (26),
\[
\|g_{K}\|\le D^{-1}e_{K}\le \frac{D}{c_{\star}}\frac{1}{K}\xrightarrow[K\to\infty]{}0.
\]
Thus $\|\nabla f(x_{k})\|\to0$, establishing part (a) of the theorem.

All steps above are justified under the stated assumptions; therefore the
conclusions of Theorem \ref{thm:global} hold. $\blacksquare$

%%%%%%%%%%%%%%%%%%%%%%%%%%%%%%%%%%%%%%%%%%%%%%%%%%%%%%%%%%%%%%%%%%%%%%%%%%%%%%%%
\section*{Rate Derivation (With Constants)}
\addcontentsline{toc}{section}{Rate Derivation (With Constants)}
From Step 8 we obtained
\[
f(x_{K})-f^{*}\le \frac{D^{2}}{c_{\star}}\frac{1}{K},
\qquad 
c_{\star}= \frac{\eta}{2L}\wedge\frac{\eta\Delta_{\min}^{2}}{LD}.
\]
Choosing the more conservative bound $c_{\star}= \frac{\eta}{2L}$ gives the
explicit constant
\[
f(x_{K})-f^{*}\le \frac{2L D^{2}}{\eta}\,\frac{1}{K},
\tag{27}
\]
which is precisely the rate (8) in Theorem \ref{thm:global}.

%%%%%%%%%%%%%%%%%%%%%%%%%%%%%%%%%%%%%%%%%%%%%%%%%%%%%%%%%%%%%%%%%%%%%%%%%%%%%%%%
\section*{Special Cases}
\addcontentsline{toc}{section}{Special Cases}
\begin{enumerate}[label=\alph*.]
\item \textbf{Exact Gradient Descent.} If $\Delta_{k}\ge\|\nabla f(x_{k})\|/L$ for all $k$,
the algorithm always takes the full step $p_{k}=-(1/L)\nabla f(x_{k})$.
The acceptance ratio $\rho_{k}=1$ (by Lemma \ref{lem:errorbound}), so every step is
accepted and $\Delta_{k}$ never shrinks. The method reduces to classical
gradient descent with stepsize $1/L$, and the bound (27) becomes the classical
gradient‑descent rate.

\item \textbf{Adaptive Radius without Upper Bound.} If we remove the cap
$\Delta_{\max}$ and let the radius grow whenever $\rho_{k}\ge\eta_{\text{inc}}$,
the method can eventually satisfy the full‑step condition, after which the
convergence rate matches that of gradient descent. The proof above remains valid
because all inequalities use only the lower bound $\Delta_{\min}$.

\item \textbf{Strongly Convex $f$.} If in addition $f$ is $\mu$‑strongly convex,
the sequence $\{f(x_{k})-f^{*}\}$ contracts linearly:
\[
f(x_{k+1})-f^{*}\le\bigl(1-\tfrac{\eta\mu}{L}\bigr)\bigl(f(x_{k})-f^{*}\bigr).
\]
The proof follows by replacing (22) with the stronger inequality
$\|g_{k}\|\ge\mu\|x_{k}-x^{*}\|$.
\end{enumerate}

%%%%%%%%%%%%%%%%%%%%%%%%%%%%%%%%%%%%%%%%%%%%%%%%%%%%%%%%%%%%%%%%%%%%%%%%%%%%%%%%
\section*{References}
\addcontentsline{toc}{section}{References}
\begin{enumerate}
\item Nocedal, J., & Wright, S. (2006). \emph{Numerical Optimization}. Springer.
\item Conn, A. R., Gould, N. I., & Toint, P. L. (2000). \emph{Trust Region Methods}. SIAM.
\item Boyd, S., & Vandenberghe, L. (2004). \emph{Convex Optimization}. Cambridge University Press.
\end{enumerate}

\end{document}